\par\medskip
\hypertarget{native-080-article061}{}\pdfbookmark[2]{61.}{article-061}
\begin{center}61.\end{center}
Tradidimus in praecedentibus methodum facillimam, corporis coelestis locum geocentricum respectu eclipticae seu aequatoris, a parallaxi liberum sive ea affectum, ac perinde a nutatione liberum seu ea affectum determinandi. Quod enim attinet ad nutationem, omnis differentia in eo versabitur, utrum aequatoris positionem mediam adoptemus an veram, adeoque in casu priori longitudines ab aequinoctio medio, in posteriori a vero numeremus, sicuti in casu illo eclipticae obliquitas media, in hoc vera adhibenda est. Ceterum sponte elucet, quo plures abbreviationes in calculo coordinatarum introducantur, eo plures operationes praeliminares esse instituendas: quamobrem praestantia methodi supra explicatae, coordinatas immediate ex anomalia excentrica deducendi, tunc potissimum se manifestabit, ubi multos locos geocentricos determinare oportet: contra quoties unus tantum locus computandus esset, aut perpauci, neutiquam operae pretium foret, laborem tot quantitates auxiliares calculandi suscipere. Quin potius in tali casu methodum vulgarem haud deserere praestabit, secundum quam ex anomalia excentrica deducitur vera atque radius vector; hinc locus heliocentricus respectu eclipticae; hinc longitudo \SourceLeaf{081}et latitudo geocentrica, atque hinc tandem rectascensio et declinatio. Ne quid igitur hic deesse videatur, duas ultimas operationes adhuc breviter explicabimus.
